% X. Cross-domain discussion. States how the pre-registered hypothesis came
% out, from the decision rules of each results/*/PREREGISTERED.md
% (paper/decisions.py). Does not generalise from one trace, capacity or cost
% (constraint 7). Generated: Fig. F8 (Table T8 in the appendix).

\section{Cross-Domain Discussion}\label{sec:discussion}

\begin{figure*}[t]
\centering
\includegraphics[width=\textwidth]{figures/f8_utility_envelopes}
\caption{Simulated utility $U(c)$ of every policy (grey) and its upper envelope
(coloured by the winning policy) against the cost of one prefetch. Left: disk,
C1, 1\,MiB, means over \diskEnvelopeTracesOneMiB{} paired traces, with the two cost
models marked. Right: web test split, logistic ranker, capacity 8, with the
assumed costs marked. The envelope is computed in hindsight on the test data
and is not a selection.}
\label{fig:f8}
\end{figure*}

\textbf{One objective, two domains.} Fig.~\ref{fig:f8} puts both domains on
the axis of Eq.~\eqref{eq:utility}. On disk, the deepest read-ahead
(\diskEnvelopeBestZeroOneMiB) is best when a prefetch is free, window growth
(\diskEnvelopeBestAssumedOneMiB) at the assumed cost model's $c$, and no
prefetching at the break-even model's $c$ above one. On the web, the best
policy at every marked positive cost is an adaptive controller that chooses
its depth per view (\webEnvelopeBestCostZeroPointZeroOne{} at $c$ =
\webHeadlineCost, \webEnvelopeBestCostZeroPointTwoFive{} at the largest marked
cost). In both domains the policy that wins changes with the price of a
prefetch, and it is a choice of how much to fetch.

\textbf{The hypothesis, by its rules.} The pre-registered hypothesis had three
parts. \emph{The predictor is the smallest lever:} not established in either
study. Every learned predictor lost to its simple, fair baseline (the class
router to \texttt{sequential}, the depth network to the validation-chosen
depth, the logistic ranker to transition counts, the GRU to online counts),
but on disk the feedback gate moved the hit ratio as little or less, and on
the web the GRU's loss could not be separated from the depth lever's gain.
\emph{Depth and freshness are large levers:} supported by the orderings. On
disk, depth and window growth are larger than both the predictor and the gate
at both cache sizes; on the web, freshness is larger than the ranker. \emph{Feedback is a large lever:} not supported. The disk
gate is no better than \texttt{sequential} and the web gate shows no
detectable difference from the adaptive controller; both mainly issue less.
The one gated disk mode above \texttt{sequential}, \texttt{depth\_adaptive},
also chooses its depth (an exploratory comparison).

\textbf{Simulation against measurement.} The two real-system studies point in
opposite directions, and both are measured. On a hard disk with O\_DIRECT and
a closed loop, the simulator's higher hit ratio turned into higher latency in
most cells,
plausibly because prefetch reads compete with demand reads for one device.
Over HTTP, the link serves demand first, a prefetched object skips both the
round trip and the disk, and the instant-prefetch model \emph{understated}
the saving on the hard-disk store.
The simulator's hit ratio is therefore a statement about the cache, not about
latency, and which way the error goes depends on where the bottleneck sits.

\textbf{For a system designer.} Choose how much to prefetch per workload and
per price of a prefetch, prefer window growth that merges reads when device
operations are the cost, keep the statistics fresh, and validate on the real
device before trusting a simulated hit ratio. A learned predictor must beat
these baselines on held-out data to earn its training cost; none of ours did.
